\chapter{Cost and Choice}
\label{chap:cost-choice}

\section{Learning objectives}

After completing this chapter, you should be able to:
\begin{itemize}
  \item distinguish economic cost from accounting cost;
  \item formulate a firm's cost-minimisation problem;
  \item explain the relationship between isoquants and isocost lines; and
  \item formulate the firm's profit-maximisation problem.
\end{itemize}

\section{Economic cost}

The economic cost of an action is its opportunity cost. A firm that owns a building still faces a cost from using it in production: the building could be rented to another business or sold. \keyterm{Economic profit} is revenue minus all opportunity costs, including the opportunity cost of the owner's time and capital.

\keyterm{Accounting profit} deducts recorded monetary expenses from revenue. It may therefore exceed economic profit when the firm uses resources it already owns. The distinction matters when asking whether a firm should remain in an activity, not merely whether it has paid its bills.

\section{Cost minimisation}

Suppose labour costs $w$ per unit and capital costs $r$ per unit. The cost of an input bundle is
\[
C=wL+rK.
\]
For a given output target $\bar q$, the firm chooses the least-cost input combination capable of producing it:
\[
\min_{L,K}\; wL+rK
\quad\text{subject to}\quad f(L,K)\geq \bar q.
\]

The constraint says that the firm must reach the target. Under monotonic technology, cost is minimised on the isoquant $f(L,K)=\bar q$: purchasing inputs beyond that isoquant would add cost without raising the required output.

\section{Isocost lines}

An \keyterm{isocost line} shows all input bundles with the same total cost. Solving $C=wL+rK$ for $K$ gives
\[
K=\frac{C}{r}-\frac{w}{r}L.
\]
Its slope is $-w/r$, the negative of the relative price of labour in terms of capital. Higher isocost lines represent higher expenditure.

With smooth technology, the least-cost bundle is the point where an isoquant is tangent to the lowest attainable isocost line. The tangency condition is
\[
MRTS_{L,K}=\frac{MP_L}{MP_K}=\frac{w}{r}.
\]
The firm substitutes between inputs until its technological trade-off equals the market trade-off implied by input prices.

\section{Changes in input prices}

If the wage $w$ rises while $r$ is fixed, isocost lines become steeper. Labour becomes relatively more expensive, so a cost-minimising firm generally substitutes capital for labour when the technology allows it. The extent of this response depends on the elasticity of substitution between inputs.

With Leontief technology, substitution is impossible: the firm must use fixed proportions regardless of relative prices. With perfect-substitute inputs, the firm uses only the cheaper input per unit of effective service.

\section{Revenue and profit}

Revenue is the money received from sales. If output is $q$ and price is $p$, total revenue is $R(q)=p(q)q$. Profit is
\[
\pi(q)=R(q)-C(q).
\]
The firm's problem is to choose output, and indirectly inputs, to maximise profit. The precise form of the problem depends on market structure. A competitive firm treats price as given; a monopolist recognises that selling more output changes market price.

The general marginal principle is that output should be expanded while marginal revenue is greater than marginal cost. At an interior profit maximum,
\[
MR(q)=MC(q).
\]
This condition will be developed in the following chapter and applied to competitive and monopolistic markets later in the book.

\section{Short-run and long-run choice}

In the short run, fixed inputs create costs that cannot be avoided immediately. In the long run, the firm can adjust all inputs, enter a market, exit it, or change scale. A firm can therefore make a loss in the short run yet continue operating if revenue covers variable costs and operating reduces the loss relative to shutting down. The distinction between fixed and variable cost is developed next.

\section*{Chapter Review}

\subsection*{Key concepts}
\begin{itemize}
  \item \textbf{Economic cost}: opportunity cost of all resources used.
  \item \textbf{Isocost line}: input bundles with the same total cost.
  \item \textbf{Cost minimisation}: producing a target output at the lowest feasible cost.
  \item \textbf{Economic profit}: total revenue minus economic cost.
  \item \textbf{Profit maximisation}: choosing output where marginal revenue equals marginal cost, subject to relevant conditions.
\end{itemize}

\section*{Exercises}

\subsection*{Conceptual questions}
\begin{enumerate}[label=\arabic*.]
  \item Distinguish accounting profit from economic profit. Give an example of an implicit cost.
  \item Define an isocost line and interpret its slope.
  \item Explain why a cost-minimising firm normally chooses a point on, rather than above, the target isoquant.
  \item Explain how a rise in wages affects a cost-minimising firm's input choice.
\end{enumerate}

\subsection*{Graphical exercises}
\begin{enumerate}[resume,label=\arabic*.]
  \item Draw an isoquant and three parallel isocost lines. Mark the cost-minimising input bundle.
  \item Show how an increase in the wage changes the slope of an isocost line.
  \item Draw a Leontief isoquant and explain why the tangency condition is not applicable at the optimum.
\end{enumerate}

\subsection*{Analytical exercises}
\begin{enumerate}[resume,label=\arabic*.]
  \item A firm uses $q=L^{1/2}K^{1/2}$, with $w=4$ and $r=1$. Write the cost-minimisation condition for producing a given output.
  \item Let $R(q)=80q-2q^2$ and $C(q)=10+8q$. Find $MR$, $MC$, and the profit-maximising quantity.
  \item A shop owner uses a premises she owns. The premises could be rented out for \euro2,000 per month. Explain how this affects economic and accounting profit.
\end{enumerate}

\subsection*{Exam-style multiple-choice questions}
\begin{enumerate}[resume,label=\arabic*.]
  \item The slope of an isocost line with labour on the horizontal axis and capital on the vertical axis is: (a) $-w/r$; (b) $-r/w$; (c) $w+r$; or (d) $C/r$.
  \item At a smooth interior cost minimum: (a) $MRTS=w/r$; (b) $MRTS=r/w$; (c) $MR=MC$; or (d) $AP=MP$.
  \item Economic profit differs from accounting profit because economic profit: (a) ignores revenue; (b) includes opportunity costs; (c) excludes wages; or (d) must be zero.
\end{enumerate}

\subsection*{PPE discussion question}

\noindent Firms frequently argue that a regulation raises their costs. Discuss which costs are relevant for public policy: firms' accounting costs, opportunity costs, costs borne by consumers, and costs imposed on third parties.
